```latex
\documentclass[11pt]{article}
\usepackage{amsmath,amssymb,amsthm}

\newtheorem{theorem}{Theorem}
\newtheorem{lemma}{Lemma}

\begin{document}

\section{Answer}

No. The statement is false in general. We give a counterexample with \(P=\mathbb N\).

\section{Counterexample}

Let \(P=\mathbb N\), the additive monoid of nonnegative integers. It is integral, saturated, and sharp.

\begin{lemma}
For \(P=\mathbb N\):
\begin{enumerate}
\item \(P[-0]^\sharp\simeq P\), with structure map \(\mathrm{id}_P\colon P\to P\).
\item For every \(n>0\), \(P[-n]^\sharp\simeq 0\), with structure map the unique map \(P\to 0\).
\end{enumerate}
\end{lemma}

\begin{proof}
For \(x=0\), the condition that \(x\) is sent to \(0\) is vacuous, and the initial object among sharp monoids under \(P\) is \(P\xrightarrow{\mathrm{id}_P}P\).

For \(x=n>0\), let \(f\colon P\to Q\) be a homomorphism of monoids with \(f(n)=0\). Then \(n f(1)=f(n)=0\) in \(Q\). Since \(Q\) is integral and sharp, this forces \(f(1)=0\). Hence \(f\) is the zero homomorphism \(P\to 0\). Thus the initial such homomorphism is \(P\to 0\).
\end{proof}

Let
\[
a=(P\xrightarrow{\mathrm{id}_P}P),\qquad b=(P\to 0)
\]
be the corresponding objects of \(A_P\). By the lemma, \(B_P\) has, up to isomorphism, the two objects \(a\) and \(b\), and there is a unique morphism \(a\to b\) in \(A_P\). Thus in \(B_P^{\mathrm{op}}\) there is a unique non-identity morphism \(b\to a\).

Now observe that in \(A_P\), the object \(b=(P\to 0)\) is terminal: for every object \(X=(P\to Q)\), the unique monoid homomorphism \(Q\to 0\) gives the unique \(A_P\)-morphism \(X\to b\). Therefore, in \(A_P^{\mathrm{op}}\), the object \(b\) is initial.

Let \(F\) be the presheaf of abelian groups on \(B_P^{\mathrm{op}}\) defined by
\[
F(a)=\mathbb Z,\qquad F(b)=0,
\]
with the unique restriction map \(F(a)\to F(b)\) equal to the zero map. This is a well-defined presheaf, and since \(B_P^{\mathrm{op}}\) has the trivial topology, \(F\) is a sheaf.

Let \(i\colon B_P^{\mathrm{op}}\to A_P^{\mathrm{op}}\) be the inclusion. For any object \(X\in A_P^{\mathrm{op}}\), consider the comma category
\[
(i\downarrow X)=\{(Y,\alpha)\mid Y\in B_P^{\mathrm{op}},\ \alpha\colon i(Y)\to X\}.
\]
Since \(b\) is initial in \(A_P^{\mathrm{op}}\) and belongs to \(B_P^{\mathrm{op}}\), the unique morphism \(!\colon b\to X\) gives an object \((b,!)\) of \((i\downarrow X)\). For any other object \((Y,\alpha)\), the unique morphism \(\gamma\colon b\to Y\) in \(A_P^{\mathrm{op}}\) lies in \(B_P^{\mathrm{op}}\) because the subcategory is full, and
\[
\alpha\circ i(\gamma)=!
\]
by uniqueness of morphisms out of the initial object \(b\). Hence \((b,!)\) is an initial object of \((i\downarrow X)\).

The left Kan extension \(i_!F\) is given by the standard colimit formula
\[
(i_!F)(X)=\operatorname*{colim}_{(Y,\alpha)\in (i\downarrow X)^{\mathrm{op}}} F(Y).
\]
Since \((b,!)\) is initial in \((i\downarrow X)\), it is terminal in \((i\downarrow X)^{\mathrm{op}}\), and therefore the colimit is
\[
(i_!F)(X)=F(b)=0.
\]
Thus \(i_!F=0\) as a presheaf on \(A_P^{\mathrm{op}}\).

Sheafification of the zero presheaf is the zero sheaf, so
\[
\pi^{-1}F=0.
\]
Consequently,
\[
\mathrm R\pi_*\pi^{-1}F=\mathrm R\pi_*(0)=0.
\]
The natural morphism
\[
F\longrightarrow \mathrm R\pi_*\pi^{-1}F
\]
is therefore the zero morphism \(F\to 0\). But \(F\) is not quasi-isomorphic to \(0\), because its cohomology sheaf in degree \(0\) is \(F\) itself, and
\[
F(a)=\mathbb Z\neq 0.
\]
Thus the morphism is not a quasi-isomorphism.

Hence the proposed statement fails for \(P=\mathbb N\).

\end{document}
```